% =============================================================================
% 12-pipelining-risc-cisc.tex -- Week 12: instruction pipelining, RISC/CISC,
% workload-based design choice, and course review
%
% Citation discipline:
%   * Stallings 11e Ch. 17 is the indexed RISC/CISC anchor: §17.5 RISC
%     pipelining (PDF p. 730) and §17.9 CISC/RISC (PDF p. 752).
%   * Stallings 11e Ch. 20 §20.1 is used only for the brief stream taxonomy
%     connection (PDF p. 847); the chapter is not treated as a pipeline text.
%   * P&H ARM 5e Ch. 6 (pp. 516--579) supports workload/parallelism framing;
%     App. D (pp. D-2--D-43) supports the RISC survey and instruction-format
%     comparison. P&H Ch. 4 is used for the standard five-stage model.
%
% Compile from Slides/CS401 with Windows MiKTeX/XeLaTeX:
%   TEXINPUTS="../../Preambles;" xelatex -interaction=nonstopmode 12-pipelining-risc-cisc.tex
%   BIBINPUTS="../..;" bibtex 12-pipelining-risc-cisc
%   TEXINPUTS="../../Preambles;" xelatex -interaction=nonstopmode 12-pipelining-risc-cisc.tex
%   TEXINPUTS="../../Preambles;" xelatex -interaction=nonstopmode 12-pipelining-risc-cisc.tex
% =============================================================================

\documentclass[aspectratio=169]{beamer}
\input{header}
\coursecode{CS401}
\pdfstringdefDisableCommands{\def\\{, }}

\title{Instruction Pipelining and RISC vs. CISC}
\subtitle{Week 12 --- Throughput, Workload Fit, and the Course's Final Trade-offs}
\author[Dr. Auqib Hamid Lone]{%
  \textbf{Dr. Auqib Hamid Lone}\\
  \small Assistant Professor in Computer Science \& Engineering
}
\institute[GCET Ganderbal]{%
  \small\textbf{PCC CS-401: Computer Organization and Architecture}\\[0.3em]
  Department of Computer Science \& Engineering\\
  Government College of Engineering \& Technology (GCET), Ganderbal
}
\date[\today]{\small\today}

\begin{document}

\frame{\titlepage}

% =============================================================================
% ACT 0 --- THE COURSE ARC
% =============================================================================
\section{The Course Arc}

\begin{frame}{Where Week 11 Left Us}
  \begin{block}{Overlap is useful only when meaning is preserved}
    \begin{itemize}
      \item Week 11 classified parallel work and diagnosed structural, data,
        and control hazards in a five-stage instruction stream.
      \item Forwarding, stalls, scheduling, prediction, and recovery each
        exchange hardware or control complexity for useful throughput.
      \item The question for the final week is broader: \key{which instruction
        and processor choices make that overlap easier to sustain?}
    \end{itemize}
  \end{block}
  \begin{exampleblock}{Today's bridge}
    Pipelining is the mechanism; RISC and CISC are different philosophies for
    exposing and implementing the instruction work that the pipeline must move.
  \end{exampleblock}
\end{frame}

\begin{frame}{Motivation: A Fast Clock Is Not a Fast Program}
  \begin{exampleblock}{Two processors, one workload}
    Processor A completes 1 instruction per cycle at 2~GHz. Processor B has a
    3~GHz clock, but its average CPI is 2. What should we compare?
  \end{exampleblock}
  \begin{itemize}
    \item The performance equation from Week 1 still governs the answer:
      \[
        T_{CPU}=\frac{\mathrm{IC}\times\mathrm{CPI}}{f}.
      \]
    \item A pipeline can reduce the \emph{effective} CPI for a long stream,
      but fills, stalls, flushes, cache misses, and instruction count remain.
    \item ISA, compiler, datapath, memory hierarchy, and workload meet in the
      final measured time. No single label wins in isolation.
  \end{itemize}
  \muted{The performance framing follows the course's Week 1 model and P\&H
  Ch.~1, §1.6 (p.~28) \cite{PattersonHennessy2017_computer_organization_design}.}
\end{frame}

\begin{frame}{Roadmap: Four Questions to Close CS401}
  \begin{enumerate}
    \item How do latency, throughput, fill, drain, and hazards determine
      pipeline speedup?
    \item What does a RISC philosophy make regular, and what does a CISC
      philosophy make expressive?
    \item How should a workload decide between competing design priorities?
    \item Can we turn the whole course into a small set of reusable questions
      for the review and doubt-clearing session?
  \end{enumerate}
  \vspace{0.15cm}
  \begin{alertblock}{Running comparison}
    The same operation, $\texttt{A[i] = A[i] + 1}$, will be viewed as an
    instruction stream, a pipeline workload, and an ISA-design choice.
  \end{alertblock}
\end{frame}

\transitionslide{First quantify what the pipeline actually improves.}

% =============================================================================
% ACT 1 --- PIPELINE QUANTITATIVE MODEL
% =============================================================================
\section{Pipelining: Latency and Throughput}

\begin{frame}{The Five-Stage Model, Revisited}
  \begin{block}{One instruction stream, five stage functions}
    \begin{tabular}{@{}lllll@{}}
      \textbf{IF} & \textbf{ID} & \textbf{EX} & \textbf{MEM} & \textbf{WB}\\
      Fetch & Decode/read & ALU/branch & Data memory & Register write
    \end{tabular}
  \end{block}
  \begin{itemize}
    \item Pipeline registers separate stages and carry the instruction's
      control and data state across clock boundaries.
    \item The clock period is constrained by the slowest stage plus pipeline
      register overhead.
    \item The instruction's \key{latency} is the time from its IF to its WB;
      the stream's \key{throughput} is the interval between completed results.
  \end{itemize}
  \muted{The standard five-stage datapath and hazard vocabulary are indexed in
  P\&H Ch.~4, §§4.5--4.8 (pp.~283, 297, 316, 328).}
\end{frame}

\begin{frame}{Worked Example 1: Fill, Drain, and Speedup}
  \scriptsize
  \begin{exampleblock}{Assumptions}
    $n=6$ independent instructions, $k=5$ balanced stages, one cycle per
    stage, no hazards, and no extra pipeline-register cost.
  \end{exampleblock}
  \begin{center}
    \begin{tabular}{@{}c|cccccccccc@{}}
      \toprule
      & 1 & 2 & 3 & 4 & 5 & 6 & 7 & 8 & 9 & 10\\
      \midrule
      $I_1$ & IF & ID & EX & MEM & WB & & & & &\\
      $I_2$ & & IF & ID & EX & MEM & WB & & & &\\
      $I_3$ & & & IF & ID & EX & MEM & WB & & &\\
      $I_4$ & & & & IF & ID & EX & MEM & WB & &\\
      $I_5$ & & & & & IF & ID & EX & MEM & WB &\\
      $I_6$ & & & & & & IF & ID & EX & MEM & WB\\
      \bottomrule
    \end{tabular}
  \end{center}
  \begin{itemize}
    \item Sequential time: $T_{seq}=nk=6\times5=30$ cycles.
    \item Pipelined time: $T_{pipe}=k+(n-1)=5+5=10$ cycles.
    \item Speedup: $S=30/10=3.0\times$, below the ideal $5\times$ because
      this finite stream pays one fill and one drain.
  \end{itemize}
\end{frame}

\begin{frame}{Latency Versus Throughput: Do Not Swap Them}
  \begin{block}{The common-case statement}
    After a $k$-stage pipeline fills, an ideal stream can complete one
    instruction every cycle. That does \emph{not} make each instruction's
    latency one cycle.
  \end{block}
  \begin{itemize}
    \item Ideal latency: approximately $kT_c$ for one instruction.
    \item Ideal steady-state throughput: one completed instruction per $T_c$.
    \item For $n$ independent instructions:
      \[
        T_{pipe}=\bigl(k+n-1\bigr)T_c,\qquad
        S_{ideal}=\frac{nk}{k+n-1}.
      \]
    \item As $n\to\infty$, $S_{ideal}\to k$ only if stages are balanced and
      the stream supplies enough independent work.
  \end{itemize}
\end{frame}

\begin{frame}{Worked Example 2: A Pipeline with Bubbles}
  \scriptsize
  \begin{exampleblock}{Load-use dependence}
    Assume the load's data is available after MEM. The immediately following
    ALU instruction needs it at EX, so forwarding still leaves one bubble.
  \end{exampleblock}
  \begin{center}
    \begin{tabular}{@{}c|ccccccccc@{}}
      \toprule
      & 1 & 2 & 3 & 4 & 5 & 6 & 7 & 8 & 9\\
      \midrule
      \texttt{LW R1,0(R2)} & IF & ID & EX & MEM & WB & & & &\\
      \texttt{ADD R3,R1,R4} & & IF & ID & \textbf{stall} & EX & MEM & WB & &\\
      \texttt{SUB R5,R3,R6} & & & IF & \textbf{stall} & ID & EX & MEM & WB &\\
      \texttt{OR R7,R8,R9} & & & & & IF & ID & EX & MEM & WB\\
      \bottomrule
    \end{tabular}
  \end{center}
  \begin{itemize}
    \item Without the bubble, the ADD would request the value before MEM has
      produced it: a true RAW hazard.
    \item With one bubble, the load result can be forwarded to ADD's EX input;
      the stream completes in 9 rather than the ideal 8 cycles.
    \item Achieved CPI is therefore shaped by the instruction mix and hazard
      rate, not by the nominal number of stages alone.
  \end{itemize}
\end{frame}

\begin{frame}{Socratic Check: What Improved?}
  \begin{alertblock}{Question}
    A five-stage pipeline completes six instructions in ten cycles. Did every
    instruction become five times faster? What evidence would you report?
  \end{alertblock}
  \begin{itemize}
    \item \bad{No:} each instruction still traverses five stages in the model.
    \item \good{The stream improved:} total time fell from 30 to 10 cycles,
      giving a $3\times$ speedup for this finite workload.
    \item Report latency and throughput separately; then list bubbles, flushes,
      memory stalls, and fill/drain effects that explain the gap from ideal.
  \end{itemize}
\end{frame}

\begin{frame}{Pipeline Design Knobs}
  \begin{tabular}{@{}p{2.8cm}p{3.3cm}p{3.3cm}p{2.8cm}@{}}
    \toprule
    \textbf{Choice} & \textbf{Potential gain} & \textbf{New cost} & \textbf{Question}\\
    \midrule
    Deeper pipeline & Shorter stage delay; higher $f$ & More registers and branch penalty & Is useful work predictable?\\
    Duplicate resources & Fewer structural hazards & Area, power, verification & Is the conflict frequent?\\
    Forwarding & Fewer RAW stalls & Muxes, comparators, timing paths & Does the value arrive in time?\\
    Prediction & Fewer control bubbles & Predictor state and flush recovery & What is $P_{flush}$?\\
    \bottomrule
  \end{tabular}
  \vspace{0.2cm}
  \muted{Stallings 11e treats instruction pipelining in §16.4 (PDF
  pp.~667--685); its Ch.~17 RISC discussion connects regular instruction
  execution to pipelining (see §17.5, PDF p.~730).}
\end{frame}

\transitionslide{Now ask what the instruction set asks the pipeline to do.}

% =============================================================================
% ACT 2 --- RISC/CISC
% =============================================================================
\section{RISC and CISC}

\begin{frame}{RISC and CISC Are Design Philosophies}
  \begin{block}{The useful contrast}
    \begin{itemize}
      \item \key{RISC:} a relatively small, regular instruction set with
        simple operations and a register-centered execution model.
      \item \key{CISC:} a richer instruction set in which individual
        instructions may express more complex operations or addressing forms.
    \end{itemize}
  \end{block}
  \begin{itemize}
    \item These are tendencies, not claims that every RISC instruction is
      simple or every CISC instruction is slow.
    \item The real comparison is between the complete implementation:
      encoding, decode, micro-operations, compiler, memory system, and target
      workload.
  \end{itemize}
  \muted{Stallings 11e Ch.~17 covers RISC execution, RISC pipelining, and the
  RISC/CISC contemporary-systems comparison in §17.9 (PDF p.~752).}
\end{frame}

\begin{frame}{Running Comparison: \texttt{A[i] = A[i] + 1}}
  \small
  \begin{exampleblock}{The operation is the same; the instruction stream differs}
    The processor must load an element, add one, and store the result. The ISA
    decides where those operands live and how much work one instruction names.
  \end{exampleblock}
  \begin{tabular}{@{}p{5.6cm}p{5.6cm}@{}}
    \toprule
    \textbf{Register-centered / RISC tendency} & \textbf{Memory-rich / CISC tendency}\\
    \midrule
    \texttt{LOAD R1,[R2+R3]} & \texttt{INC [R2+R3]}\\
    \texttt{ADD R1,R1,1} & \textit{one architectural instruction may name the memory update}\\
    \texttt{STORE [R2+R3],R1} & \textit{decode/execute may internally split the work}\\
    More instructions, regular formats, explicit loads/stores & Fewer visible instructions, richer formats, more decode work\\
    \bottomrule
  \end{tabular}
  \vspace{0.15cm}
  \muted{This is a conceptual comparison, not a claim about one specific ISA's
  exact encoding. P\&H App.~D surveys RISC addressing modes and formats
  (pp.~D-2--D-43).}
\end{frame}

\begin{frame}{What Regularity Buys a Pipeline}
  \begin{itemize}
    \item Fixed or clustered instruction formats simplify fetch boundaries and
      decode timing.
    \item A large register file reduces frequent memory operands and exposes
      values to forwarding paths; this is a RISC design tendency, not a law.
    \item Load/store separation makes the memory stage's role predictable:
      ALU instructions can use EX without also requesting data memory.
    \item Predictable stage work makes balancing and hazard detection easier,
      while the compiler carries more scheduling responsibility.
  \end{itemize}
  \begin{alertblock}{But regularity is not free}
    More instructions can increase $\mathrm{IC}$, register traffic, code size,
    instruction-cache pressure, or branch exposure for a particular workload.
  \end{alertblock}
\end{frame}

\begin{frame}{What Expressiveness Buys a CISC Design}
  \begin{itemize}
    \item A complex instruction can reduce the visible instruction count for
      a dense operation or a legacy software interface.
    \item Rich addressing modes can express array, pointer, stack, or string
      operations compactly.
    \item Compatibility and code density can matter when instruction memory,
      binary size, or an established software ecosystem dominates.
    \item Internally, a modern implementation may translate complex
      instructions into simpler micro-operations before scheduling them.
  \end{itemize}
  \begin{block}{The cost moves rather than disappears}
    Decode width, variable-length boundaries, micro-operation expansion,
    precise exceptions, and verification can complicate the front end.
  \end{block}
  \muted{Stallings Ch.~17 presents RISC and CISC as contemporary implementation
  choices, not as a binary ``fast versus slow'' ranking.}
\end{frame}

\begin{frame}{Socratic Check: Which Statement Is Too Strong?}
  \begin{alertblock}{Choose the statement that needs correction}
    (A) RISC often favors regular formats and register operands.\\
    (B) CISC can express more work per architectural instruction.\\
    (C) RISC always has fewer instructions and CISC always has lower CPI.\\
  \end{alertblock}
  \begin{itemize}
    \item (C) is too strong: instruction count, CPI, clock rate, stalls, code
      size, and workload interact through the performance equation.
    \item A fair comparison fixes the workload and measures total execution
      time, not one attractive ISA statistic.
  \end{itemize}
\end{frame}

\begin{frame}{Worked Example 3: Workload-Based Design Choice}
  \scriptsize
  \begin{exampleblock}{Two workload sketches}
    \textbf{Embedded control:} tight power budget, predictable loops, limited
    memory, fixed sensor operations.\\
    \textbf{Cloud analytics:} large software base, diverse branches, high
    throughput, frequent vector/multicore parallelism.
  \end{exampleblock}
  \begin{tabular}{@{}p{3.0cm}p{4.0cm}p{4.0cm}@{}}
    \toprule
    \textbf{Criterion} & \textbf{Embedded control} & \textbf{Cloud analytics}\\
    \midrule
    Likely priority & Predictability, energy, compact implementation & Throughput, compatibility, scalable parallel work\\
    Helpful tendency & Regular pipeline, simple decode, compiler-visible registers & Strong ecosystem, capable front end, SIMD/vector and multicore support\\
    Main risk & Code expansion and instruction-memory pressure & Decode/recovery complexity and memory bottlenecks\\
    \bottomrule
  \end{tabular}
  \vspace{0.15cm}
  \begin{itemize}
    \item Neither row proves a universal winner; it identifies what to measure.
    \item P\&H Ch.~6 emphasizes that parallel performance depends on the
      program and its communication/coordination costs (pp.~516--579).
  \end{itemize}
\end{frame}

\begin{frame}{Worked Example 4: Compare Total Time, Not Labels}
  \begin{exampleblock}{Same task, different measured parameters}
    A RISC-like implementation executes $1.20\times10^9$ instructions at
    CPI $=1.1$ and $f=2.4$~GHz. A CISC-like implementation executes
    $0.80\times10^9$ instructions at CPI $=2.0$ and $f=3.0$~GHz.
  \end{exampleblock}
  \begin{align*}
    T_R &= \frac{1.20\times10^9\times1.1}{2.4\times10^9}
         = 0.55~\text{s},\\
    T_C &= \frac{0.80\times10^9\times2.0}{3.0\times10^9}
         \approx 0.533~\text{s}.
  \end{align*}
  \begin{itemize}
    \item The CISC-like system wins this workload by about $3.0\%$ despite
      its higher CPI, because it executes fewer instructions at a higher clock.
    \item Change the workload, cache behavior, or branch rate and the result
      can change. The measurement is the conclusion; the labels are hypotheses.
  \end{itemize}
\end{frame}

\begin{frame}{Pipelining Meets the ISA Boundary}
  \begin{tabular}{@{}p{3.1cm}p{4.2cm}p{4.2cm}@{}}
    \toprule
    \textbf{Pipeline concern} & \textbf{Regular instruction stream helps} & \textbf{Rich instruction stream may require}\\
    \midrule
    Fetch/decode & Predictable boundaries and widths & Boundary detection and wider decode\\
    Data hazards & Explicit register dependencies & Internal micro-op dependencies\\
    Structural hazards & Regular stage/resource demand & Variable resource demand per instruction\\
    Control recovery & Explicit branch behavior & Precise recovery across expanded work\\
    \bottomrule
  \end{tabular}
  \vspace{0.2cm}
  \begin{block}{Course-level connection}
    Week 5's addressing modes, Weeks 6--7's control units, Week 9--10's memory
    hierarchy, and Week 11's hazards all reappear here as constraints on
    instruction throughput.
  \end{block}
\end{frame}

\transitionslide{The final design question is always: fit to which workload?}

% =============================================================================
% ACT 3 --- COURSE REVIEW AND DOUBT CLEARING
% =============================================================================
\section{Course Review}

\begin{frame}{A Single Performance Thread Through CS401}
  \begin{enumerate}
    \item \key{Representation:} bits, two's complement, and IEEE 754 determine
      what values the hardware can encode.
    \item \key{Execution:} registers, buses, addressing modes, and control
      units move and transform those values.
    \item \key{Communication:} I/O and interrupts connect the processor to
      devices without making every transfer a busy wait.
    \item \key{Storage:} caches, virtual memory, and secondary storage trade
      capacity, latency, protection, and persistence.
    \item \key{Overlap:} parallel organizations and pipelines raise throughput
      when hazards are identified and paid for explicitly.
    \item \key{Choice:} the ISA and organization should be judged by the
      workload's measured time, energy, reliability, and software constraints.
  \end{enumerate}
\end{frame}

\begin{frame}{Doubt-Clearing: The Six Questions to Ask}
  \begin{block}{When a problem looks unfamiliar}
    \begin{enumerate}
      \item What is the representation: bits, fields, signedness, or format?
      \item Where is the operand: register, memory, cache, page, or device?
      \item Which hardware state changes: PC, IR, MAR/MDR, register, flag, or
        control word?
      \item What is the bottleneck: clock, CPI, bus, memory, I/O, or hazard?
      \item What is the unit of comparison: latency, throughput, speedup,
        hit rate, miss penalty, or total CPU time?
      \item Which assumption makes the shortcut valid, and what breaks it?
    \end{enumerate}
  \end{block}
  \muted{A strong exam solution names the model before doing arithmetic.}
\end{frame}

\begin{frame}{Socratic Check: Final Integrator}
  \begin{alertblock}{Scenario}
    A workload has many independent vector operations, frequent branches, and
    a large working set. It runs on a multicore processor with a pipelined
    front end. Which earlier course topics should you inspect first?
  \end{alertblock}
  \begin{itemize}
    \item \key{Week 11 / Week 12:} SIMD or MIMD opportunity, pipeline hazards,
      branch penalty, and instruction-stream choice.
    \item \key{Weeks 9--10:} cache locality, TLB/page behavior, and storage
      misses that can dominate an otherwise fast pipeline.
    \item \key{Week 1:} total time through $\mathrm{IC}\times\mathrm{CPI}/f$;
      never infer performance from clock rate alone.
  \end{itemize}
\end{frame}

\begin{frame}{Exam and Assignment Review Checklist}
  \begin{itemize}
    \item \textbf{Performance:} write the equation, normalize units, and
      compare total time or speedup.
    \item \textbf{Instruction cycle/control:} name the state and write RTL
      transfers with \texttt{PC}, \texttt{IR}, \texttt{MAR}, and \texttt{MDR}.
    \item \textbf{Memory:} split addresses into tag/index/offset or
      page/offset before translating or computing a hit.
    \item \textbf{I/O:} identify who waits, who owns the bus, and where the
      data moves.
    \item \textbf{Pipeline:} draw a cycle table; classify each hazard; count
      every stall and flush.
    \item \textbf{RISC/CISC:} state the workload, measured variables, and the
      trade-off before defending a design choice.
  \end{itemize}
\end{frame}

\begin{frame}{Closing the Arc: From Bits to Design Judgment}
  \footnotesize
  \begin{block}{The course's recurring question}
    How can a finite amount of hardware keep the common work moving quickly,
    correctly, and predictably?
  \end{block}
  \begin{itemize}
    \item Representation makes operations possible.
    \item Organization moves operands and generates control.
    \item Hierarchy hides slow capacity behind fast common cases.
    \item Parallelism overlaps work, while hazards expose the dependencies.
    \item RISC/CISC choices shape the instruction work; the workload decides
      whether the trade-off pays.
  \end{itemize}
  \begin{exampleblock}{Final takeaway}
    \key{Do not ask which architecture is universally best. Ask which design
    makes this workload's useful work per unit time highest, and show the
    assumptions that make your answer true.}
  \end{exampleblock}
\end{frame}

\begin{frame}{References and Source Boundaries}
  \small
  \begin{itemize}
    \item Stallings, \emph{Computer Organization and Architecture: Designing
      for Performance}, 11th ed. (repository key retained as Stallings2015):
      Ch.~16 §16.4, ``Instruction Pipelining'' (PDF pp.~667--685); Ch.~17,
      especially §17.5 (PDF p.~730) and §17.9 (PDF p.~752); Ch.~20 §20.1,
      ``Multiple Processor Organizations'' (PDF p.~847).
    \item Patterson and Hennessy, \emph{Computer Organization and Design:
      The Hardware/Software Interface}, ARM 5e: Ch.~4 §§4.5--4.8
      (pp.~283, 297, 316, 328); Ch.~6 §§6.2--6.3 (pp.~518, 523); App.~D
      (pp.~D-2--D-43).
    \item The worked comparisons in this deck are constructed examples. Their
      numbers illustrate the performance equation and pipeline model; they are
      not quotations or measurements from either book.
  \end{itemize}
\end{frame}

\begin{frame}[allowframebreaks]{Bibliography}
  \bibliographystyle{plain}
  \bibliography{Bibliography_base}
\end{frame}

\end{document}